\section{Verification of the exact final machine}
\label{sec:verification}

\subsection{The universal interpreter theorem}

The selected table was checked afresh against the independently calibrated
increment/drain source. First, the checker reconstructs and compares the
complete ordered source control graph, including both outcomes of every
decrement and all break, continue and halt successors. The eleven later
closed reassociations and the two boundary rotations change addresses;
they do not change this graph. The separate closed-loop theorem connects
that graph to the original K230 source.

The tape invariant describes a marker, a binary source address, a zero
gap, a header, nine primary unary registers, a first spare register and
an arbitrary finite spare tail. Each register value is independent and
unbounded. The proof variables include the additional gap and the spare
values; they are not fixed to the small values used in diagnostic runs.
Each zero or positive decrement guard has its own proof case.

Four generic register-file operations are justified by explicit
finite-list induction schemas. Unary scans are justified by natural-number
induction on their run length. Every use of these schemas is derived
against the exact supplied transition table. The remaining transitions
are executed literally on symbolic segments. For every nonhalting source instruction, the endpoint comparison checks
the entire tape representation, head position, interpreter state, updated
registers and source successor, including the zero exterior. The halt
case reaches the unique terminal state.

The complete fresh result is:
\begin{center}
\begin{tabularx}{\linewidth}{@{}lX@{}}\toprule
Obligation & Replayed result\\\midrule
Source leaves & $5\,121$\\
Guarded instruction cases, including halt & $6\,656$\\
Symbolic rule applications & $1\,615\,284$\\
Largest single proof & 412 symbolic rules\\
Minimum actual instruction duration & 13 literal transitions\\
Native working states / transition rows & $226/452$\\
Rows outside runtime and initialization & none\\
Blank initialization & $409\,835$ transitions\\
Initialization tape support & 351 one-cells, every other cell zero\\
Source maximum depth / zero run & $41/6$\\
\bottomrule\end{tabularx}
\end{center}
The proof-construction fuel is a limit on searching for a symbolic proof.
It does not bound a register, a list, the actual duration of an instruction,
or the running time of the complete machine. All generated proof cases, the four block schemas, row coverage and exact initialization support
are retained under \path{evidence/runtime226/}.

\subsection{Initialization and state-count accounting}

The blank initialization is a separate finite calculation, because the
runtime invariant must first be established. In inherited interpreter
coordinates the head reaches $-5$, the marker is at $-6$, and the register
header is at 176. The 181 cells from the head at $-5$ through coordinate 175 are zero. The physical
primary values are
\[
                    (0,0,0,3,0,0,3,0,0),
\]
followed by 167 spare values equal to one. The physical-to-logical register
permutation produces the exact seed used by the numerical proof enumerator.
The independent calculation uses the complete set of one positions, not
just a bounded window or a register summary.

The complete final table has two rows for every working state, no row for
the separate halt label, no other undefined target, and only left/right
moves. Every counted state is reachable in the directed transition graph
from the start; this is state-count accounting and does not assert that
every state occurs in the actual blank execution. The dense table is
bijectively related to the original-label and inherited-label versions.
All transition identities and start/halt mappings are checked. The
canonicalization produced internally by the symbolic driver is a further
renaming, not a different candidate; the public table is the one with
SHA--256 beginning \path{b58279b8}.

\subsection{Composition into the consistency theorem}

\begin{theorem}[The selected machine]
Let $M_{226}$ be the literal table \path{machines/k226.tm}, started in
state zero on a blank two-sided binary tape. Then
\[
\mathrm{ZFC}\vdash
\left(\Halts(M_{226},\mathrm{blank})\leftrightarrow
          \exists p\,\mathrm{Proof}_{\mathrm{ZFC}}(p,\bot)\right).
\]
Consequently ZFC proves its nonhalting equivalent to \(\Con(\mathrm{ZFC})\).
\end{theorem}
\begin{proof}
The exact initialization enters the interpreter invariant at the
calibrated seed. The universal instruction theorem simulates every source
instruction by a finite positive number of native transitions and
reflects the unique halt instruction. Induction therefore equates native
halting with source halting. Positive progress ensures that an infinite
source execution cannot be compressed into a finite native run.

The complete source graph audit transports the final source to the
increment/drain source. Each of its four changed closed loops has the
same terminating register transformation as the original loop; contextual
replacement therefore transports halting to the pinned K230 source.
The reference enumerator and finite-derivation translations identify
that halting assertion with a classical ZFC derivation of falsity, as
spelled out in the calibration section.

Every new finite table check and every new induction here has a direct
arithmetization in ZFC. The inherited calibration also consists of
explicit syntax, proof enumeration and finite derivation translations.
Formalizing those constructions gives the displayed internal assertion.
Classical negation gives the nonhalting version. No consistency or
soundness assumption is used to prove either biconditional.
\end{proof}

\subsection{External Rocq verification}
\input{sections/rocq_status.tex}

The external theorem and the assertion of ZFC-provability are different
claims. Even a closed Rocq theorem does not by itself constitute a
serialized ZFC derivation of the arithmetized sentence. The latter is
justified here by the explicit finitary internalization argument. The
package identifies exactly which new and inherited files were compiled, which endpoint assumptions were printed, and which independent kernel
rechecks completed.

\subsection{Diagnostic checks and rejected controls}

Finite native tests were used to discover and reject candidates. The final
226 table passed 26,624 instruction fixtures and the blank initialization,
covering $15\,414\,531$ literal transitions in that test family. These
tests are useful diagnostics, but the universal theorem above supplies
the arbitrary-register claim.

The rejected seven-zero trials are especially informative: a source
reassociation can preserve the register program and still break the
binary decoder. Its failing address, source branch and tape discrepancy
are retained with a replay script. The two repairs at the 227-state stage and the five repairs at the
226-state stage are therefore necessary parts of the construction.

A bounded further native search on the repaired 227 table found no
additional plain or one/two-helper fold with radius five and local blocks
of at most 65 steps. The recorded search covered 25,651 state pairs,
46,377,008 row choices and 585,927,699 helper checks. This is evidence about
that specified search family only. It is not a lower bound on possible
machines or on other compilation schemes.
